\subsection{What the Abel spectral-phase limit requires}
\label{dev:spectralphases}

\paragraph{The recorded question and a defensible target.}
Occurrence 1975-391.7 and problem GP-1975-24 concern the expression on
\GSpage{P000652}. Let \(M\) be a closed compact Riemannian manifold,
let \(\lambda_k\geq0\) be the eigenvalues of \(\sqrt{-\Delta}\), counted
with multiplicity, and fix \(a>0\). With the periodic delta normalized to
have mass one on \([0,2\pi)\), set
\[
 \Phi_{a,\varepsilon}=\varepsilon\sum_k e^{-\varepsilon\lambda_k}
 \left(\delta_{\rm per}(\mu-a\lambda_k)-\frac1{2\pi}\right),
 \qquad \varepsilon>0.
\]
The notebook proposes a limit and relates its nonzero Fourier coefficients
to resonances with lengths of closed geodesics. It also specifies a generic
length arrangement. To make that phrase coherent, distinctness and
incommensurability must concern primitive unoriented geodesics; iterates
necessarily have commensurable lengths. Isolated, nondegenerate iterates
are the modern hypotheses needed for the simple leading wave-trace
singularity. They are not silently imposed on every manifold in the first
line of the notes. Analyticity, as recorded there, can be retained but
does not itself prove the limiting assertion.

The strongest target justified here is conditional: identify every fixed
Fourier coefficient by a local trace calculation, then prove a bound
uniform in both the Fourier index and the Abel parameter. The latter
step cannot be replaced by pointwise convergence of coefficients. The
following development isolates this obstruction, solves the complete
circle model, and checks what fails when the generic-geodesic restriction
is dropped. No general theorem under the notebook's generic hypotheses
is claimed.

\paragraph{Literature and identification limit.}
The leading singularity at a nondegenerate orbit is part of the
wave-trace theory already published in 1975
\citep{duistermaatguillemin1975spectrum}.
The first expansion in \citet[Introduction]{zelditch1998wave} explicitly
has a simple-pole leading term and logarithmic lower terms. That local
expansion supports the fixed-mode mechanism; it does not, by itself,
give a bound uniform over all sampled times \(ma\).
The latter paper appeared in 1998 and was deposited on arXiv in 2000;
the deposit is not its first-publication date.
As of 10 November 2026, the searches recorded for this chapter did not
identify a theorem establishing the entire notebook limit in the stated
generality. This is an unresolved match and analytic gap in the argument
presented here, not a claim that all spectral-phase limit problems remain
open in the literature.

\begin{lemma}[A necessary and sufficient Fourier test]\ours\quad
For \(T_\varepsilon\in\mathcal D'(\mathbb R/2\pi\mathbb Z)\),
\(0<\varepsilon\leq1\), write
\(\widehat T_\varepsilon(m)=(2\pi)^{-1}
\langle T_\varepsilon,e^{-im\mu}\rangle\).
Suppose the family is bounded on each compact parameter interval
away from zero. It converges in distributions as \(\varepsilon\downarrow0\)
if and only if every coefficient converges and there are
\(C,N,\varepsilon_0>0\) such that
\[
 |\widehat T_\varepsilon(m)|\leq C(1+|m|)^N
 \quad(m\in\mathbb Z,\ 0<\varepsilon\leq\varepsilon_0).
\]
For the notebook family, \(\widehat\Phi_{a,\varepsilon}(0)=0\), and
\[
 \widehat\Phi_{a,\varepsilon}(m)
 =\frac{\varepsilon}{2\pi}
    \sum_k e^{-(\varepsilon+ima)\lambda_k},\qquad m\ne0.
\]
\end{lemma}
\begin{proof}
The last formula follows by pairing each delta with the indicated
exponential; the subtracted constant removes only the zero mode.
For necessity, distributional convergence gives pointwise boundedness
on the Frechet space of smooth periodic test functions. The uniform
boundedness principle gives equicontinuity: for some \(C,N\),
\(|\langle T_\varepsilon,f\rangle|\leq
C\max_{j\leq N}\|f^{(j)}\|_\infty\) for small \(\varepsilon\).
Testing exponentials gives the polynomial bound and their individual
convergence. For sufficiency, the limiting coefficients have the same
polynomial bound and so define a distribution. The Fourier coefficients
of a smooth test function decay faster than any power. Consequently
\[
 \langle T_\varepsilon,f\rangle
   =2\pi\sum_{m\in\mathbb Z}
       \widehat T_\varepsilon(m)\widehat f(-m)
\]
converges by dominated convergence. This proves the assertion and keeps
the normalization explicit.
\end{proof}

\begin{proposition}[The complete circle model]\ours\quad
Let \(M\) be the circle of circumference \(2\pi R\), \(R>0\).
If \(a/(2\pi R)\) is irrational, then
\(\Phi_{a,\varepsilon}\to0\) in distributions. If
\(a/(2\pi R)=p/q\) in lowest terms with \(p,q\) positive integers, then
\[
 \Phi_{a,\varepsilon}\longrightarrow
 2R\left(\frac1q\sum_{j=0}^{q-1}
           \delta_{\rm per}(\mu-2\pi j/q)-\frac1{2\pi}\right).
\]
Including or omitting the zero eigenvalue does not change the limit.
\end{proposition}
\begin{proof}
The eigenfunctions \(e^{ikx/R}\), \(k\in\mathbb Z\), give frequencies
\(|k|/R\). Thus, for \(m\ne0\), the geometric series gives
\[
 \widehat\Phi_{a,\varepsilon}(m)
 =\frac{\varepsilon}{2\pi}
 \left(1+\frac{2}{e^{(\varepsilon+ima)/R}-1}\right).
\]
Its limit is \(R/\pi\) when \(ma/R\in2\pi\mathbb Z\), and zero
otherwise. Moreover,
\[
 \left|\widehat\Phi_{a,\varepsilon}(m)\right|
 \leq\frac{\varepsilon}{2\pi}
       \left(1+\frac{2}{e^{\varepsilon/R}-1}\right)
 \leq\frac{\varepsilon+2R}{2\pi},
\]
uniformly in \(m\). The Fourier test therefore proves convergence.
In the rational case, the resonant indices are precisely the nonzero
multiples of \(q\), which are the coefficients of the displayed atomic
comb. In the irrational case no nonzero index is resonant. A zero
eigenvalue contributes only
\(\varepsilon(\delta_{\rm per}-1/(2\pi))\), which vanishes.
\end{proof}

\begin{example}[Why one cannot omit the geometric restriction]\ours\quad
On the unit round two-sphere, with \(a=2\pi\), the notebook family has
no distributional limit. More precisely,
\[
 \widehat\Phi_{2\pi,\varepsilon}(1)
     =-\frac1{\pi\varepsilon}+O(1).
\]
This example does \emph{not} refute the generic-geodesic assertion:
the round sphere has continuous families of closed geodesics of the
same length, and their return maps are degenerate.
\end{example}
\begin{proof}
The sphere frequencies are \(\lambda_\ell=\sqrt{\ell(\ell+1)}\), with
multiplicity \(2\ell+1\), \(\ell\geq0\). For completeness, the polar
coordinate identity
\(\Delta_{\mathbb R^3}=\partial_r^2+2r^{-1}\partial_r+r^{-2}\Delta_{S^2}\)
applied to a harmonic homogeneous polynomial \(r^\ell Y\) gives
\(-\Delta_{S^2}Y=\ell(\ell+1)Y\).
The decomposition of degree-\(\ell\) homogeneous polynomials into
harmonic polynomials plus \(r^2\) times degree-\(\ell-2\) polynomials
gives dimension
\(\binom{\ell+2}{2}-\binom{\ell}{2}=2\ell+1\).
Polynomial restrictions are dense on the sphere, so these harmonics
give the full spectrum.

Put \(\delta_\ell=\ell+\tfrac12-\lambda_\ell\).
For \(\ell\geq1\), rationalization gives
\[
 0<\delta_\ell=
 \frac{1/4}{\ell+1/2+\sqrt{\ell(\ell+1)}}
       \leq\frac1{4\ell}.
\]
Uniformly for \(0<\varepsilon\leq1\),
\[
 e^{-(\varepsilon+2\pi i)\lambda_\ell}
 =-e^{-\varepsilon(\ell+1/2)}
       \bigl(1+O(\ell^{-1})\bigr).
\]
After multiplying by \(2\ell+1\), summing, and then multiplying by
\(\varepsilon/(2\pi)\), the error is \(O(1)\).
The finitely many omitted low terms do not change this bound. The main
term is
\[
 -\frac{\varepsilon}{2\pi}e^{-\varepsilon/2}
       \sum_{\ell\geq0}(2\ell+1)e^{-\varepsilon\ell}
 =-\frac{\varepsilon}{2\pi}e^{-\varepsilon/2}
       \frac{1+e^{-\varepsilon}}{(1-e^{-\varepsilon})^2}
 =-\frac1{\pi\varepsilon}+O(1).
\]
A distributional limit would have a finite first Fourier coefficient,
which this estimate excludes.
\end{proof}

\paragraph{Checks, scope, and next step.}
The circle normalization was checked independently by summing over
residue classes \(k\bmod q\): each class has Abel mass tending to
\(R/q\), doubled by the two signs of nonzero \(k\). Its Fourier
transform agrees with the geometric-series calculation. The uniform
bound follows separately from the triangle inequality for the positive
Abel sum. For the sphere, rationalization gives a uniform remainder;
the closed sum of \((2\ell+1)e^{-\varepsilon\ell}\) verifies the sign
and the coefficient of the divergence. A separate numerical check,
\texttt{code/spectral\_phase\_check.py}, compares the circle formula
with direct summation and the sphere asymptotic with successively
smaller parameters; the output is in
\texttt{computations/spectral\_phase\_check.json}. Those floating-point
values are diagnostics, not the proof of convergence or divergence.

The remaining question in the generic higher-dimensional setting is
the uniform polynomial estimate in the Fourier test, together with
the exact hypotheses under which the local trace coefficients are
well defined and do not cancel. A useful further research run would
trace Shnirelman's exact publication and establish such a bound for
one specified geometric class and a specified set of admissible
\(a\). The present calculation neither proves that global estimate
nor labels the sphere counterexample as a failure under hypotheses
it does not satisfy.
